% Problem Statement

Anime production has a \textbf{contribution invisibility problem}: much of the creative
labour that produces it goes unrecorded, unattributed and unverifiable.
% ----- 3.1 The Credit Gap -----
\subsection{The Credit Gap}
\label{subsec:credit_gap}

End credits are a curated subset rather than a log. Three things narrow them.
Work outsourced through successive tiers surfaces, if at all, as the
subcontracting company's name, so the individuals at the second or third tier
appear nowhere. Screen time and production-committee policy cap how many names
run at all. And an animator who leaves mid-production may be dropped despite
having finished cuts. \Cref{fig:cut_record} shows where this leaves the record:
the work is recorded in detail by the party that commissioned it, and the version
that reaches the public names a company.

\begin{figure*}[t]
  \centering
  \includegraphics[width=\textwidth]{figures/cut_record}
  \caption{The asymmetry this paper addresses. Work commissioned down a chain of
  subcontractors produces two records of who did it. The order the studio accepts
  names the person and the cuts it paid for, and stays inside the studio. The
  credits that reach the audience name a company. A co-signed certificate is a
  form of the first record that can cross the boundary.}
  \label{fig:cut_record}
\end{figure*}

That a remedy is wanted is not ours to assert, and we do not have to: it appears in the
industry's own survey. Asked what should change, a production assistant answering
JAniCA's 2023 survey proposed a licence issued inside production---by the in-between
checker for in-betweening, by the character designer or chief animation director for key
animation---so that a creator's record travels with
them~\cite{janica2023survey}.\footnote{Free-text answer at p.~102, and individual
testimony rather than a survey statistic. We cite it because the proposal is structurally
the one this paper develops, and because it was made from inside production rather than
by a vendor. In the 2026 wave a finishing inspector puts the complaint in one line: the
work is done and it is not credited~\cite{janica2026survey}.}

% ----- 3.2 Consequences for Freelance Animators -----
\subsection{Consequences for Freelance Animators}
\label{subsec:consequences}

The credit gap bears on the share of animators who work outside employment
(\cref{subsec:labor}), and what it costs them follows from what a credit is for.
What fails here is a market for a good whose quality the buyer cannot observe
before purchase, which is the problem Akerlof
described~\cite{akerlof1970lemons} and to which Spence's answer was a signal
costly enough that only those who can bear it will send
it~\cite{spence1973signaling}. Creative labour markets run on exactly this
substitution. Menger's review of the field finds that in many art worlds initial
training is an imperfect filtering device, so what sorts practitioners is not a
credential~\cite{menger1999artistic}, and in project-based production careers are
assembled from a sequence of engagements rather than from tenure, so each
engagement has to be legible to the next~\cite{faulkner1987projects}. A credit is
the signal such a market produces. Withholding it does not merely deny
recognition; it removes the instrument the market runs on.

This framing also sets the bar the rest of the paper has to clear. A signal
separates types only if it is costly to send falsely. We return to whether a
studio countersignature meets that condition in \cref{subsec:trust}.

The consequences are ordinary and compounding. A freelancer seeking the next contract has
to demonstrate a track record, and without credits that means verbal claims, personal
networks or samples nobody can independently attribute, while the studio hiring them has
no standard way to check. The same absence removes the evidence from fee negotiation: a
key animator who delivered thirty cuts on a well-known production, uncredited, has nothing
to point at. And because the record lives in colleagues' memories, it degrades when
studios close, which they do, or when the people who remember move on.

A stronger claim is sometimes made for records of this kind, and we are not making it.
Mortgage applications, visa procedures and credit assessments require documented
\emph{income}, evidenced in Japan by tax returns and withholding statements. A
contribution record attests to work performed, not earnings, and no institution recognises
an on-chain record for these purposes. We claim only that the absence of any verifiable
professional record is one component of animators' weak institutional standing.

% ----- 3.3 Limitations of Existing Solutions -----
\subsection{Limitations of Existing Solutions}
\label{subsec:existing_limitations}

None of the responses in place produces what is missing.
\Cref{tab:provenance} classifies them by where their record comes from; here we
note only what each leaves open. The Freelance Protection Act regulates written
terms and payment, not attribution, so a contract can be fully compliant and the
animator still uncredited. NAFCA's template agreements improve bilateral
transparency but produce a document between two parties, not something the
animator can carry forward. Fan-maintained databases and portfolio platforms
record at the level the credential needs but carry no attestation, the first
resting on volunteer effort and the second on self-report.

Production-management tools are the interesting case, because they already hold
the right data. Trackers used in animation and visual effects---Autodesk Flow,
ftrack, the open-source Kitsu, and Japanese anime-specific tools---capture
per-cut assignment and completion as a byproduct of running the work. That record
is studio-owned: it lives in the studio's instance, is not an attestation the
freelancer holds, and stops being reachable at offboarding or closure. They solve
production tracking, which is a different problem from a portable credential, and
they are also the system in which an accepted order is recorded---which is why the
design in \cref{sec:framework} starts there rather than from a new data source.

% ----- 3.4 Problem Formalization -----
\subsection{Problem Formalization}
\label{subsec:formalization}

Let $A$ be the set of animators who contributed to a production $P$, and
$C(P) \subseteq A$ the set credited in its official records. The \textbf{credit gap} is:

\begin{equation}
  \Delta(P) = A \setminus C(P)
  \label{eq:credit_gap}
\end{equation}

For any animator $a \in \Delta(P)$ there is no verifiable, persistent, portable record of
their contribution to $P$. The definition is conceptual rather than a computed quantity,
since the uncredited set is by construction unrecorded. A practical proxy would compare,
for a sample of productions, the contributors a production-management tool records against
the names in those productions' broadcast credits; we propose that as future work and do
not estimate $|\Delta(P)|$ here.

An ideal solution should satisfy the following requirements:

\begin{enumerate}
  \item[\textbf{R1}] \textbf{Completeness:} Every animator $a \in A$ can obtain a record,
    regardless of their position in the subcontracting hierarchy.
  \item[\textbf{R2}] \textbf{Verifiability:} Third parties can cryptographically verify
    the authenticity of the record without contacting the original studio.
  \item[\textbf{R3}] \textbf{Persistence:} The record survives studio closures,
    platform shutdowns, and changes in industry infrastructure.
  \item[\textbf{R4}] \textbf{Portability:} The animator owns and controls the record
    and can present it to any future employer or client.
  \item[\textbf{R5}] \textbf{Minimal friction:} The recording process must integrate
    into existing workflows without requiring additional effort from the animator.
\end{enumerate}

Read against these requirements, each mechanism of
\cref{subsec:existing_limitations} fails a different one, and none fails only R5.

AnimatorSBT is designed to satisfy all five requirements, with one important
qualification recorded here and revisited in \cref{subsec:requirements_eval}:
for \textbf{R2 (Verifiability)} it currently provides only \emph{integrity and
existence} verification (that a claim was recorded, unaltered, at a given time),
not \emph{authenticity} (that the claimed work was actually performed).
Authenticity requires the studio co-signing layer we design and implement as a
reference contract in \cref{subsec:cosigning} but do not yet deploy on chain or
operate with an independent studio. R4 and R5 are satisfied by design but rest on
wallet abstraction that is not built, and have not been validated with users.
